\frfilename{mt274.tex}
\versiondate{13.4.10}
\copyrightdate{1995}

\def\chaptername{Teori probabilitas}
\def\sectionname{Teorema limit pusat}

\newsection{274}

Teorema besar kedua yang menjadi pokok bab ini mempunyai jenis yang baru.
Teorema ini merupakan teorema limit, tetapi limit yang terlibat adalah limit
{\it distribusi}, bukan limit fungsi (seperti dalam teorema limit kuat di atas
atau teorema martingal di bawah), dan bukan pula limit kelas ekuivalensi fungsi
(seperti dalam Bab 24). Saya memberikan tiga bentuk teorema ini, dalam
274I-274K, semuanya diturunkan sebagai akibat Teorema 274G; pembuktiannya
tersebar sepanjang 274C-274G. Dalam 274A-274B dan 274M saya memberikan
sifat-sifat paling dasar dari distribusi normal.

\leader{274A}{Distribusi normal }\cmmnt{Kita memerlukan beberapa fakta dari
teori probabilitas dasar.

\medskip

}{\bf (a)}\cmmnt{ Ingat bahwa

\Centerline{$\int_{-\infty}^{\infty}e^{-x^2/2}dx=\sqrt{2\pi}$}

\noindent (263G). Oleh karena itu, jika kita menetapkan}\dvro{ Jika kita menetapkan}{}

\Centerline{$\mu_G E=\Bover1{\sqrt{2\pi}}\int_Ee^{-x^2/2}dx$}

\noindent untuk setiap himpunan terukur Lebesgue $E$, maka $\mu_G$ adalah
ukuran probabilitas Radon\cmmnt{ (256E)}; kita menyebutnya {\bf distribusi
normal baku}. Fungsi distribusi yang bersesuaian adalah

\Centerline{$\Phi(a)
\cmmnt{\mskip5mu =\mu_G\ocint{-\infty,a}}
=\Bover1{\sqrt{2\pi}}\int_{-\infty}^ae^{-x^2/2}dx$}

\noindent untuk $a\in\Bbb R$; di seluruh bagian ini saya akan mencadangkan
lambang $\Phi$ untuk fungsi tersebut.

\cmmnt{Dengan menulis $\Sigma$ untuk aljabar himpunan bagian terukur
Lebesgue dari $\Bbb R$, $(\Bbb R,\Sigma,\mu_G)$ adalah ruang probabilitas.
Perhatikan bahwa ruang ini lengkap, dan mempunyai himpunan-himpunan terabaikan
yang sama dengan ukuran Lebesgue, karena $e^{-x^2/2}>0$ untuk setiap
$x$ (bandingkan 234Lc).}

\header{274Ab}{\bf (b)} Variabel acak $X$ disebut {\bf normal baku}
jika distribusinya adalah $\mu_G$\cmmnt{; yakni, jika fungsi
$x\mapsto\Bover1{\sqrt{2\pi}}e^{-x^2/2}$ merupakan fungsi kepadatan bagi $X$.
Maksud dari catatan dalam (a) adalah bahwa variabel acak demikian memang ada;
misalnya, ambil ruang probabilitas $(\Bbb R,\Sigma,\mu_G)$ di sana, dan
tetapkan $X(x)=x$ untuk setiap $x\in\Bbb R$}.

\header{274Ac}{\bf (c)} Jika $X$ adalah variabel acak normal baku, maka

\Centerline{$\Expn(X)
=\Bover1{\sqrt{2\pi}}\int_{-\infty}^{\infty}xe^{-x^2/2}dx
=0$,}

\Centerline{$\Var(X)
=\Bover1{\sqrt{2\pi}}\int_{-\infty}^{\infty}x^2e^{-x^2/2}dx
=1$\dvro{.}{}}

\cmmnt{\noindent berdasarkan 263H.}

\header{274Ad}{\bf (d)} Secara lebih umum, variabel acak $X$ disebut {\bf
normal} jika terdapat $a\in\Bbb R$ dan $\sigma>0$ sedemikian sehingga
$Z=(X-a)/\sigma$ normal baku. Dalam kasus ini\cmmnt{ $X=\sigma Z+a$ sehingga}
$\Expn(X)\cmmnt{\mskip5mu =\sigma\Expn(Z)+a}=a$,
$\Var(X)\cmmnt{\mskip5mu =\sigma^2\Var(Z)}=\sigma^2$.
\cmmnt{

Untuk setiap $c\in\Bbb R$, kita mempunyai

$$\eqalignno{\Bover1{\sigma\sqrt{2\pi}}
  \int_{-\infty}^ce^{-(x-a)^2/2\sigma^2}dx
&=\Bover1{\sqrt{2\pi}}\int_{-\infty}^{(c-a)/\sigma}
  e^{-y^2/2}dy\cr
\noalign{\noindent (dengan menyubstitusikan $x=a+\sigma y$ untuk
$-\infty<y\le(c-a)/\sigma$)}
&=\Pr(Z\le\Bover{c-a}{\sigma})
=\Pr(X\le c).\cr}$$

\noindent Jadi}
$x\mapsto\Bover1{\sigma\sqrt{2\pi}}e^{-(x-a)^2/2\sigma^2}$
adalah fungsi kepadatan bagi $X$\cmmnt{ (271Ib)}.
Sebaliknya,\cmmnt{ tentu saja,} variabel acak
dengan fungsi kepadatan seperti ini adalah normal,
dengan ekspektasi $a$ dan varians $\sigma^2$.
{\bf Distribusi normal} adalah distribusi dengan fungsi-fungsi kepadatan ini.

\header{274Ae}{\bf (e)} Jika\cmmnt{ $Z$ normal baku, maka $-Z$ juga normal
baku, karena

\Centerline{$\Pr(-Z\le a)=\Pr(Z\ge -a)
=\Bover1{\sqrt{2\pi}}\int_{-a}^{\infty}e^{-x^2/2}dx
=\Bover1{\sqrt{2\pi}}\int_{-\infty}^{a}e^{-x^2/2}dx$.}

\noindent Definisi dalam kalimat pertama (d) sekarang memperlihatkan dengan
jelas bahwa jika} $X$ normal, maka $a+bX$ juga normal untuk setiap
$a\in\Bbb R$ dan $b\in\Bbb R\setminus\{0\}$.

\leader{274B}{Proposisi} Misalkan $X_1,\ldots,X_n$ adalah variabel acak
normal yang saling bebas. Maka $Y=X_1+\ldots+X_n$ normal, dengan
$\Expn(Y)=\Expn(X_1)+\ldots+\Expn(X_n)$ dan
$\Var(Y)=\Var(X_1)+\ldots+\Var(X_n)$.

\proof{ Ada tak terhitung banyaknya bukti untuk fakta ini; bukti berikut
memberi saya kesempatan untuk memperlihatkan kekuatan Bab 26, tetapi tentu
saja (dengan harga berupa aljabar yang agak tidak menyenangkan) 272U juga
memberikan hasil tersebut.

\medskip

{\bf (a)} Tinjau dahulu kasus $n=2$. Dengan menetapkan
$a_i=\Expn(X_i)$, $\sigma_i=\sqrt{\Var(X_i)}$,
$Z_i=(X_i-a_i)/\sigma_i$
kita memperoleh variabel normal baku $Z_1$, $Z_2$ yang saling bebas. Tetapkan
$\rho=\sqrt{\sigma_1^2+\sigma_2^2}$, dan nyatakan $\sigma_1$, $\sigma_2$
sebagai $\rho\cos\theta$, $\rho\sin\theta$. Tinjau $U=\cos\theta
Z_1+\sin\theta Z_2$. Kita tahu bahwa $(Z_1,Z_2)$ mempunyai fungsi kepadatan

\Centerline{$(\zeta_1,\zeta_2)\mapsto g(\zeta_1,\zeta_2)=
\Bover1{2\pi}e^{-(\zeta_1^2+\zeta_2^2)/2}$}

\noindent (272I). Oleh karena itu, untuk setiap $c\in\Bbb R$,

\Centerline{$\Pr(U\le c)=\int_Fg(z)dz$,}

\noindent dengan
$F=\{(\zeta_1,\zeta_2):\zeta_1\cos\theta+\zeta_2\sin\theta\le c\}$.
Namun sekarang misalkan $T$ adalah matriks

$$\Matrix{\cos\theta&-\sin\theta\\\sin\theta&\cos\theta}.$$

\noindent Maka mudah diperiksa bahwa

\Centerline{$T^{-1}[F]=\{(\eta_1,\eta_2):\eta_1\le c\}$,}

\Centerline{$\det T=1$,
\quad$g(Ty)=g(y)$ untuk setiap $y\in\BbbR^2$,}

\noindent sehingga berdasarkan 263A

\Centerline{$\Pr(U\le c)=\int_Fg(z)dz=\int_{T^{-1}[F]}g(Ty)dy
=\int_{\ocint{-\infty,c}\times\Bbb R}g(y)dy
=\Pr(Z_1\le c)=\Phi(c)$.}

\noindent Karena ini benar untuk setiap $c\in\Bbb R$, $U$ juga normal baku
(saya kembali menggunakan 271Ga). Tetapi

\Centerline{$X_1+X_2=\sigma_1Z_1+\sigma_2Z_2+a_1+a_2
=\rho U+a_1+a_2$,}

\noindent sehingga $X_1+X_2$ normal.

\medskip

{\bf (b)} Sekarang kita dapat melakukan induksi pada $n$. Jika $n=1$
hasilnya seketika. Untuk langkah induksi menuju $n+1\ge 2$, kita tahu bahwa
$X_1+\ldots+X_n$ normal, berdasarkan hipotesis induksi, dan bahwa $X_{n+1}$
saling bebas dengan $X_1+\ldots+X_n$, berdasarkan 272L. Jadi
$X_1+\ldots+X_n+X_{n+1}$ normal, berdasarkan (a).

Perhitungan ekspektasi dan varians $X_1+\ldots+X_n$ langsung diperoleh dari
271Ab dan 272S.
}%end of proof of 274B

\leader{274C}{Lemma} Misalkan $U_0,\ldots,U_n,V_0,\ldots,V_n$ adalah
variabel acak bernilai real yang saling bebas dan
$h:\Bbb R\to\Bbb R$ fungsi terukur Borel yang terbatas. Maka

\Centerline{$|\Expn\bigl(h(\sum_{i=0}^nU_i)-h(\sum_{i=0}^nV_i)\bigr)|
\le\sum_{i=0}^n\sup_{t\in\Bbb R}|\Expn\bigl(h(t+U_i)-h(t+V_i)\bigr)|$.}

\proof{ Untuk $0\le j\le n+1$, tetapkan
$Z_j=\sum_{i=0}^{j-1}U_i+\sum_{i=j}^nV_i$, dengan mengambil
$Z_0=\sum_{i=0}^nV_i$ dan $Z_{n+1}=\sum_{i=0}^nU_i$, dan untuk $j\le n$
tetapkan
$W_j=\sum_{i=0}^{j-1}U_i+\sum_{i=j+1}^nV_i$, sehingga $Z_j=W_j+V_j$ dan
$Z_{j+1}=W_j+U_j$ serta $W_j$, $U_j$, dan $V_j$ saling bebas (saya
menggunakan 272K, seperti dalam 272L). Maka

$$\eqalign{|\Expn\bigl(h(\sum_{i=0}^nU_i)-h(\sum_{i=0}^nV_i)\bigr)|
&=|\Expn\bigl(\sum_{i=0}^nh(Z_{i+1})-h(Z_i)\bigr)|\cr
&\le\sum_{i=0}^n|\Expn\bigl(h(Z_{i+1})-h(Z_i)\bigr)|\cr
&=\sum_{i=0}^n|\Expn\bigl(h(W_i+U_i)-h(W_i+V_i)\bigr)|.\cr}$$

Untuk menaksir jumlah ini, saya mengubahnya menjadi jumlah integral sebagai
berikut. Untuk setiap $i$, misalkan $\nu_{W_i}$ adalah distribusi $W_i$,
dan demikian seterusnya. Karena $(w,u)\mapsto w+u$ kontinu, sehingga terukur
Borel, $(w,u)\mapsto h(w+u)$ juga terukur Borel; oleh karena itu
$(w,u,v)\mapsto h(w+u)-h(w+v)$ terukur untuk setiap ukuran produk
$\nu_{W_i}\times\nu_{U_i}\times\nu_{V_i}$ pada $\BbbR^3$, dan
271E serta 272G memberi kita

$$\eqalignno{\bigl|\Expn\bigl(h(W_i+U_i)-&h(W_i+V_i)\bigr)\bigr|\cr
&=\bigl|\int h(w+u)-h(w+v)
(\nu_{W_i}\times\nu_{U_i}\times\nu_{V_i})d(w,u,v)\bigr|\cr
&=\bigl|\int\bigl(\int h(w+u)-h(w+v)(\nu_{U_i}\times\nu_{V_i})d(u,v)
\bigr)\nu_{W_i}(dw)\bigr|\cr
&\le\int\bigl|\int h(w+u)-h(w+v)(\nu_{U_i}\times\nu_{V_i})d(u,v)\bigr|
\nu_{W_i}(dw)\cr
&=\int\bigl|\Expn\bigl(h(w+U_i)-h(w+V_i)\bigr)\bigr|\nu_{W_i}(dw)\cr
&\le\sup_{t\in\Bbb R}
    \bigl|\Expn\bigl(h(t+U_i)-h(t+V_i)\bigr)\bigr|.\cr}$$

\noindent Jadi kita memperoleh

$$\eqalign{|\Expn\bigl(h(\sum_{i=0}^nU_i)-h(\sum_{i=0}^nV_i)\bigr)|
&\le\sum_{i=0}^n|\Expn\bigl(h(W_i+U_i)-h(W_i+V_i)\bigr)|\cr
&\le\sum_{i=0}^n
\sup_{t\in\Bbb R}|\Expn\bigl(h(t+U_i)-h(t+V_i)\bigr)|,\cr}$$

\noindent sebagaimana diperlukan.
}%end of proof of 274C

\leader{274D}{Lemma} Misalkan $h:\Bbb R\to\Bbb R$ adalah fungsi terbatas
yang terdiferensialkan tiga kali sedemikian sehingga
$M_2=\sup_{x\in\Bbb R}|h''(x)|$, $M_3=\sup_{x\in\Bbb R}|h'''(x)|$
keduanya berhingga. Misalkan $\epsilon>0$.

(a) Misalkan $U$ adalah variabel acak bernilai real dengan ekspektasi nol dan
varians berhingga $\sigma^2$. Maka untuk setiap $t\in\Bbb R$ kita mempunyai

\Centerline{$|\Expn(h(t+U))-h(t)-\Bover{\sigma^2}2h''(t)|
\le\Bover16\epsilon M_3\sigma^2+M_2\Expn(\psi_{\epsilon}(U))$}

\noindent dengan $\psi_{\epsilon}(x)=0$ jika $|x|\le \epsilon$, dan $x^2$
jika $|x|>\epsilon$.

(b) Misalkan $U_0,\ldots,U_n,V_0,\ldots,V_n$ adalah variabel acak yang
saling bebas dengan varians berhingga, dan anggap bahwa
$\Expn(U_i)=\Expn(V_i)=0$ serta
$\Var(U_i)=\Var(V_i)=\sigma_i^2$ untuk setiap $i\le n$. Maka

$$\eqalign{|\Expn\bigl(h(\sum_{i=0}^n&U_i)-h(\sum_{i=0}^nV_i)\bigr)|\cr
&\le\Bover13\epsilon M_3\sum_{i=0}^n\sigma_i^2
+M_2\sum_{i=0}^n\Expn\bigl(\psi_{\epsilon}(U_i)\bigr)
+M_2\sum_{i=0}^n\Expn\bigl(\psi_{\epsilon}(V_i)\bigr).\cr}$$

\proof{{\bf (a)} Intinya adalah bahwa, berdasarkan teorema Taylor dengan
sisa,

\Centerline{$|h(t+x)-h(t)-xh'(t)|\le\Bover12M_2x^2$,}

\Centerline{$|h(t+x)-h(t)-xh'(t)-\Bover12x^2h''(t)|
\le\Bover16M_3|x|^3$}

\noindent untuk setiap $x\in\Bbb R$. Jadi

\Centerline{$|h(t+x)-h(t)-xh'(t)-\Bover12x^2h''(t)|
\le\min(\Bover16M_3|x|^3,M_2x^2)
\le \Bover16\epsilon M_3x^2+M_2\psi_{\epsilon}(x)$.}

\noindent Dengan mengintegralkan terhadap distribusi $U$, kita memperoleh

$$\eqalign{|\Expn\bigl(h(t+U))-h(t)-\Bover12h''(t)\sigma^2\bigr)|
&=|\Expn(h(t+U))-h(t)-h'(t)\Expn(U)
-\Bover12h''(t)\Expn(U^2)|\cr
&=|\Expn\bigl(h(t+U)-h(t)-h'(t)U-\Bover12h''(t)U^2\bigr)|\cr
&\le\Expn\bigl(|h(t+U)-h(t)-h'(t)U-\Bover12h''(t)U^2|\bigr)\cr
&\le\Expn\bigl(\Bover16\epsilon M_3U^2+M_2\psi_{\epsilon}(U)\bigr)\cr
&=\Bover16\epsilon M_3\sigma^2+M_2\Expn(\psi_{\epsilon}(U)),\cr}$$

\noindent sebagaimana diklaim.

\medskip

{\bf (b)} Berdasarkan 274C,

$$\eqalign{|\Expn\bigl(h(\sum_{i=0}^nU_i)-h(\sum_{i=0}^nV_i)\bigr)|
&\le\sum_{i=0}^n\sup_{t\in\Bbb R}|\Expn\bigl(h(t+U_i)-h(t+V_i)\bigr)|\cr
&\le\sum_{i=0}^{n}\sup_{t\in\Bbb R}\bigl(
|\Expn(h(t+U_i))-h(t)-\Bover12h''(t)\sigma_i^2|\cr
&\qquad\qquad\qquad
+|\Expn(h(t+V_i))-h(t)-\Bover12h''(t)\sigma_i^2|\bigr),\cr}$$

\noindent yang berdasarkan (a) di atas tidak lebih dari

\Centerline{$\sum_{i=0}^n\Bover13\epsilon  M_3\sigma_i^2
+M_2\Expn(\psi_{\epsilon}(U_i))+M_2\Expn(\psi_{\epsilon}(V_i))$,}

\noindent sebagaimana diklaim.
}%end of proof of 274D

\leader{274E}{Lemma} Untuk setiap $\epsilon>0$, terdapat fungsi
$h:\Bbb R\to[0,1]$ yang terdiferensialkan tiga kali, dengan turunan ketiga
kontinu, sedemikian sehingga $h(x)=1$ untuk $x\le -\epsilon$ dan $h(x)=0$
untuk $x\ge\epsilon$.

\proof{ Misalkan $f:\ooint{-\epsilon,\epsilon}
\to\ooint{0,\infty}$
adalah sebarang fungsi yang terdiferensialkan dua kali sedemikian sehingga

\Centerline{
$\lim_{x\downarrow-\epsilon}f^{(n)}(x)
=\lim_{x\uparrow\epsilon}f^{(n)}(x)=
0$}

\noindent untuk $n=0$, $1$, dan $2$, dengan $f^{(n)}$ menyatakan turunan
ke-$n$ dari $f$; misalnya, kita dapat mengambil
$f(x)=(\epsilon^2-x^2)^3$, atau
$f(x)=\exp(-\Bover1{\epsilon^2-x^2})$. Sekarang tetapkan

\Centerline{$h(x)=1-\int_{-\epsilon}^xf/\int_{-\epsilon}^{\epsilon}f$}

\noindent untuk $|x|\le\epsilon$, lalu perluas fungsi tersebut dengan menetapkan
$h(x)=1\text{ jika }x\le-\epsilon,\quad\text{dan}\quad
h(x)=0\text{ jika }x\ge\epsilon$.
}%end of proof of 274E

\leader{274F}{Teorema Lindeberg} Misalkan $\epsilon>0$. Maka terdapat
$\delta>0$ sedemikian sehingga kapan pun $X_0,\ldots,X_n$ adalah variabel
acak bernilai real yang saling bebas sedemikian sehingga

\Centerline{$\Expn(X_i)=0$ untuk setiap $i\le n$,}

\Centerline{$\sum_{i=0}^n\Var(X_i)=1$,}

\Centerline{$\sum_{i=0}^n\Expn(\psi_{\delta}(X_i))\le\delta$}

\noindent (dengan $\psi_{\delta}(x)=0$ jika $|x|\le\delta$, dan $x^2$ jika
$|x|>\delta$), maka

\Centerline{$\bigl|\Pr(\sum_{i=0}^nX_i\le a)
-\Phi(a)\bigr|
\le\epsilon$}

\noindent untuk setiap $a\in\Bbb R$.

\proof{{\bf (a)} Misalkan $h:\Bbb R\to[0,1]$ adalah fungsi yang
terdiferensialkan tiga kali, dengan turunan ketiga kontinu, sedemikian sehingga
$\chi\ocint{-\infty,-\epsilon}\le h\le\chi\ocint{-\infty,\epsilon}$, seperti
dalam 274E. Tetapkan

\Centerline{$M_2=\sup_{x\in\Bbb
R}|h''(x)|=\sup_{|x|\le\epsilon}|h''(x)|$,}

\Centerline{$M_3=\sup_{x\in\Bbb
R}|h'''(x)|=\sup_{|x|\le\epsilon}|h'''(x)|$;}

\noindent karena \hbox{$h'''$} kontinu, keduanya berhingga. Tuliskan
$\epsilon'=\epsilon(1-\Bover2{\sqrt{2\pi}})>0$, dan ambil $\eta>0$
sedemikian sehingga

\Centerline{$(\Bover13M_3+2M_2)\eta
\le \epsilon'$.}

\noindent Perhatikan bahwa $\lim_{m\to\infty}\psi_m(x)=0$ untuk setiap $x$,
sehingga jika $X$ adalah variabel acak dengan varians berhingga, kita harus
mempunyai $\lim_{m\to\infty}\Expn(\psi_m(X))\discretionary{}{}{}=0$,
berdasarkan Teorema Konvergensi Terdominasi Lebesgue; ambil $m\ge 1$
sedemikian sehingga $\Expn(\psi_m(Z))\le\eta$, dengan $Z$ suatu (atau
sebarang) variabel acak normal baku. Akhirnya, ambil $\delta>0$ sedemikian
sehingga $\delta+\delta^2\le(\eta/m)^2$; perhatikan bahwa $\delta\le\eta$.

(Saya berharap Anda sudah cukup sering melihat argumen $\epsilon$-$\delta$
sehingga tidak terganggu oleh harapan untuk memahami alasan setiap rumus
tertentu di sini sebelum membaca seluruh argumen. Namun rumus
$\bover13M_3+2M_2$, dalam kaitannya dengan $\psi_{\delta}$, semestinya
mengingatkan kita pada 274D.)

\medskip

{\bf (b)} Misalkan $X_0,\ldots,X_n$ adalah variabel acak yang saling bebas
dengan ekspektasi nol sedemikian sehingga $\sum_{i=0}^n\Var(X_i)=1$ dan
$\sum_{i=0}^n\Expn(\psi_{\delta}(X_i))\le\delta$. Kita memerlukan barisan
bantu $Z_0,\ldots,Z_n$ dari variabel acak normal baku untuk dipasangkan dengan
$X_i$. Untuk membentuknya, saya menggunakan cara berikut. Misalkan ruang
probabilitas yang mendasari $X_0,\ldots,X_n$ adalah $(\Omega,\Sigma,\mu)$.
Tetapkan $\Omega'=\Omega\times\BbbR^{n+1}$, dan misalkan $\mu'$ adalah ukuran
produk pada $\Omega'$, dengan $\Omega$ diberi ukuran $\mu$ dan setiap faktor
$\Bbb R$ dari $\BbbR^{n+1}$ diberi ukuran $\mu_G$. Tetapkan
$X'_i(\omega,z)=X_i(\omega)$ dan $Z_i(\omega,z)=\zeta_i$ untuk
$\omega\in\dom X_i$, $z=(\zeta_0,\ldots,\zeta_n)\in\BbbR^{n+1}$,
$i\le n$. Maka $X'_0,\ldots,X'_n,Z_0,\ldots,Z_n$ saling bebas, dan setiap
$X'_i$ mempunyai distribusi yang sama dengan $X_i$ (272Mb). Akibatnya
$S'=X'_0+\ldots+X'_n$ mempunyai distribusi yang sama dengan
$S=X_0+\ldots+X_n$ (menggunakan 272T, atau dengan cara lain); sehingga
$\Expn(g(S'))=\Expn(g(S))$ untuk setiap fungsi terukur Borel terbatas $g$
(menggunakan 271E). Setiap $Z_i$ juga mempunyai distribusi $\mu_G$, sehingga
normal baku.

\medskip

{\bf (c)} Tuliskan $\sigma_i=\sqrt{\Var(X_i)}$ untuk setiap $i$, dan tetapkan
$K=\{i:i\le n,\,\sigma_i>0\}$. Perhatikan bahwa $\eta/\sigma_i\ge m$ untuk
setiap $i\in K$. \Prf\ Kita tahu bahwa

\Centerline{$\sigma_i^2=\Var(X_i)=\Expn(X_i^2)
\le\Expn(\delta^2+\psi_{\delta}(X_i))=\delta^2+\Expn(\psi_{\delta}(X_i))
\le\delta^2+\delta$,}

\noindent sehingga

\Centerline{$\Bover{\eta}{\sigma_i}\ge\Bover{\eta}{\sqrt{\delta+\delta^2}}
\ge m$}

\noindent berdasarkan pilihan $\delta$. \Qed

\medskip

{\bf (d)} Tinjau variabel acak normal $\sigma_iZ_i$ yang saling bebas.
Kita mempunyai $\Expn(\sigma_iZ_i)=\Expn(X'_i)=0$ dan
$\Var(\sigma_iZ_i)=\Var(X'_i)=\sigma_i^2$ untuk setiap $i$, sehingga
$Z=\sigma_0Z_0+\ldots+\sigma_nZ_n$ mempunyai ekspektasi $0$ dan varians
$\sum_{i=0}^n\sigma_i^2=1$; terlebih lagi, berdasarkan 274B, $Z$ normal,
sehingga sebenarnya normal baku. Sekarang kita mempunyai

$$\eqalignno{\sum_{i=0}^n\Expn(\psi_{\eta}(\sigma_iZ_i))
&=\sum_{i\in K}\Expn(\psi_{\eta}(\sigma_iZ_i))
=\sum_{i\in K}\sigma_i^2\Expn(\psi_{\eta/\sigma_i}(Z_i))\cr
\noalign{\noindent (karena
$\sigma^2\psi_{\eta/\sigma}(x)=\psi_{\eta}(\sigma x)$ kapan pun
$x\in\Bbb R$, $\sigma>0$)}
&=\sum_{i\in K}\sigma_i^2\Expn(\psi_{\eta/\sigma_i}(Z))
\le\sum_{i\in K}\sigma_i^2\Expn(\psi_{m}(Z))\cr
\noalign{\noindent (karena, berdasarkan (c),
$\eta/\sigma_i\ge m$ untuk setiap $i\in K$, sehingga
$\psi_{\eta/\sigma_i}(t)\le\psi_m(t)$ untuk setiap $t$)}
&=\Expn(\psi_m(Z))
\le\eta\cr}$$

\noindent (berdasarkan pilihan $m$). Di sisi lain, kita tentu mempunyai

\Centerline{$\sum_{i=0}^n\Expn(\psi_{\eta}(X'_i))
=\sum_{i=0}^n\Expn(\psi_{\eta}(X_i))
\le\sum_{i=0}^n\Expn(\psi_{\delta}(X_i))
\le\delta\le\eta$.}

\medskip

{\bf (e)} Untuk setiap bilangan real $t$, tetapkan

\Centerline{$h_{t}(x)=h(x-t)$}

\noindent untuk setiap $x\in\Bbb R$. Maka $h_t$ terdiferensialkan tiga kali,
dengan $\sup_{x\in\Bbb R}|h_t''(x)|=M_2$ dan
$\sup_{x\in\Bbb R}|h_t'''(x)|=M_3$. Akibatnya

\Centerline{$|\Expn(h_t(S))-\Expn(h_t(Z))|
\le\epsilon'$.}

\noindent\vthsp\Prf\ Berdasarkan 274Db,

$$\eqalign{|\Expn(h_t(S))-\Expn(h_t(Z))|
&=|\Expn(h_t(S'))-\Expn(h_t(Z))|\cr
&=|\Expn(h_t(\sum_{i=0}^nX'_i))-\Expn(h_t(\sum_{i=0}^n\sigma_iZ_i))|\cr
&\le\Bover13\eta M_3\sum_{i=0}^n\sigma_i^2
+M_2\sum_{i=0}^n\Expn(\psi_{\eta}(X_i))
+M_2\sum_{i=0}^n\Expn(\psi_{\eta}(\sigma_iZ_i))\cr
&\le\Bover13\eta M_3+M_2\eta+M_2\eta
\le\epsilon',\cr}$$

\noindent berdasarkan pilihan $\eta$. \Qed


\medskip

{\bf (f)} Sekarang ambil sebarang $a\in\Bbb R$. Kita mempunyai

\Centerline{$\chi\ocint{-\infty,a-2\epsilon}
\le h_{a-\epsilon}\le\chi\ocint{-\infty,a}
\le h_{a+\epsilon}\le\chi\ocint{-\infty,a+\epsilon}$.}

\noindent Perhatikan pula bahwa, untuk setiap $b$,

\Centerline{$\Phi(b+2\epsilon)
=\Phi(b)+\Bover1{\sqrt{2\pi}}\int_b^{b+2\epsilon}e^{-x^2/2}dx
\le\Phi(b)+\Bover{2\epsilon}{\sqrt{2\pi}}
=\Phi(b)+\epsilon-\epsilon'$.}


\noindent Akibatnya

$$\eqalign{\Phi(a)-\epsilon
&\le\Phi(a-2\epsilon)-\epsilon'
=\Pr(Z\le a-2\epsilon)-\epsilon'
\le\Expn(h_{a-\epsilon}(Z))-\epsilon'
\le\Expn(h_{a-\epsilon}(S))\cr
&\le\Pr(S\le a)\cr
&\le\Expn(h_{a+\epsilon}(S))
\le\Expn(h_{a+\epsilon}(Z))+\epsilon'
\le\Pr(Z\le a+2\epsilon)+\epsilon'
=\Phi(a+2\epsilon)+\epsilon'\cr
&\le\Phi(a)+\epsilon.\cr}$$

\noindent Namun ini tepat berarti bahwa

\Centerline{$\bigl|\Pr(\sum_{i=0}^nX_i\le a)
-\Phi(a)\bigr|
\le\epsilon$,}

\noindent sebagaimana diklaim.
}%end of proof of 274F

\leader{274G}{Teorema Limit Pusat} Misalkan $\sequencen{X_n}$ adalah
barisan variabel acak yang saling bebas, semuanya dengan ekspektasi nol dan
varians berhingga; tuliskan $s_n=\sqrt{\sum_{i=0}^n\Var(X_i)}$ untuk setiap
$n$. Anggap bahwa

\Centerline{$\lim_{n\to\infty}\Bover1{s_n^2}\sum_{i=0}^n
\Expn(\psi_{\delta s_n}(X_i))=0$ untuk setiap $\delta>0$,}

\noindent dengan $\psi_{\delta}(x)=0$ jika $|x|\le\delta$, dan $x^2$ jika
$|x|>\delta$. Tetapkan

\Centerline{$S_n=\Bover1{s_n}(X_0+\ldots+X_n)$}

\noindent untuk setiap $n\in\Bbb N$ sedemikian sehingga $s_n>0$. Maka

\Centerline{$\lim_{n\to\infty}\Pr(S_n\le a)
=\Phi(a)$}

\noindent secara seragam untuk $a\in\Bbb R$.

\proof{ Diberikan $\epsilon>0$, ambil $\delta>0$ seperti dalam teorema
Lindeberg (274F). Maka untuk semua $n$ yang cukup besar,

\Centerline{$\Bover1{s_n^2}\sum_{i=0}^n
\Expn(\psi_{\delta s_n}(X_i))\le\delta$.}

\noindent Tetapkan sebarang $n$ demikian. Tentu saja kita mempunyai $s_n>0$.
Tetapkan

\Centerline{$X'_i=\Bover1{s_n}X_i$ untuk $i\le n$;}

\noindent maka $X'_0,\ldots,X'_n$ saling bebas, dengan ekspektasi nol,

\Centerline{$\sum_{i=0}^n\Var(X'_i)=\sum_{i=0}^n\Bover1{s_n^2}\Var(X_i)
=1$,}

\Centerline{$\sum_{i=0}^n\Expn(\psi_{\delta}(X'_i))
=\sum_{i=0}^n\Bover1{s_n^2}\Expn(\psi_{\delta s_n}(X_i))
\le\delta$.}

\noindent Berdasarkan 274F,

\Centerline{$\bigl|\Pr(S_n\le a)
-\Phi(a)\bigr|
=\bigl|\Pr(\sum_{i=0}^nX'_i\le a)
-\Phi(a)\bigr|
\le\epsilon$}

\noindent untuk setiap $a\in\Bbb R$. Karena ini benar untuk semua $n$ yang
cukup besar, kita memperoleh hasilnya.
}%end of proof of 274G

\leader{274H}{Catatan (a)} Syarat

\Centerline{$\lim_{n\to\infty}\Bover1{s_n^2}\sum_{i=0}^n
\Expn(\psi_{\epsilon s_n}(X_i))=0$ untuk setiap $\epsilon>0$}

\noindent disebut {\bf syarat Lindeberg}\cmmnt{, mengikuti
{\smc Lindeberg 1922}}.

\header{274Hb}{\bf (b)} Syarat Lindeberg bersifat perlu sekaligus cukup,
dalam arti berikut. Misalkan $\sequencen{X_n}$ adalah barisan variabel acak
bernilai real yang saling bebas dengan ekspektasi nol dan varians berhingga;
tuliskan $\sigma_n=\sqrt{\Var(X_n)}$,
$s_n=\sqrt{\sum_{i=0}^n\Var(X_i)}$ untuk setiap $n$. Anggap bahwa
$\lim_{n\to\infty}s_n=\infty$,
$\lim_{n\to\infty}\Bover{\sigma_n}{s_n}=0$ dan bahwa
$\lim_{n\to\infty}\Pr(S_n\le a)=\Phi(a)$ untuk setiap $a\in\Bbb R$, dengan
$S_n=\Bover1{s_n}(X_0+\ldots+X_n)$. Maka

\Centerline{$\lim_{n\to\infty}\Bover1{s_n^2}\sum_{i=0}^n
\Expn(\psi_{\epsilon s_n}(X_i))=0$}

\noindent untuk setiap $\epsilon>0$. \cmmnt{({\smc Feller 66}, \S XV.6,
Teorema 3; {\smc Lo\`eve 77}, \S21.2.)
}%end of comment

\cmmnt{\header{274Hc}{\bf (c)}
Bukti 274F-274G di sini diadaptasi dari {\smc Feller 66}, \S VIII.4.
Keunggulannya ialah bersifat ‘elementer’, dalam arti tidak melibatkan fungsi
karakteristik. Tentu hal ini harus dibayar dengan sejumlah taksiran terperinci;
dan -- yang jauh lebih serius -- membuat kita tidak memiliki salah satu teknik
paling ampuh untuk mendeskripsikan distribusi. Bukti tersebut memang
menawarkan metode untuk membatasi

\Centerline{$|\Pr(S_n\le a)-\Phi(a)|$;}

\noindent tetapi perlu dikatakan bahwa batas yang diperoleh tidak berguna,
karena terlalu pesimistis secara mencolok, sekurang-kurangnya dalam kasus yang
mudah dianalisis. (Misalnya, dalam banyak kasus batas yang lebih baik diberikan
oleh teorema Berry-Ess\'een: jika $\sequencen{X_n}$ saling bebas dan
berdistribusi identik, dengan ekspektasi nol, dan nilai bersama dari
$\sqrt{\Expn(X_n^2)}$, $\Expn(|X_n|^3)$ masing-masing adalah
$\sigma$, $\rho<\infty$, maka

\Centerline{$|\Pr(S_n\le a)-\Phi(a)|
\le\Bover{33\rho}{4\sigma^3\sqrt{n+1}}$;}

\noindent lihat {\smc Feller 66}, \S XVI.5, {\smc Lo\`eve 77},
\S21.3, atau {\smc Hall 82}.) Selain itu, ketika $|a|$ besar,
$\Phi(a)$ amat dekat dengan $0$ atau $1$, sehingga sebarang batas seragam bagi
$|\Pr(S\le a)-\Phi(a)|$ hanya memberikan sangat sedikit informasi; banyak
sekali pekerjaan telah dilakukan untuk menaksir ekor distribusi demikian
secara lebih tepat, dengan tunduk pada syarat-syarat khusus. Misalnya, jika
$X_0,\ldots,X_n$ adalah variabel acak yang saling bebas dengan ekspektasi nol,
terbatas secara seragam dengan $|X_i|\le K$ hampir di mana-mana untuk setiap
$i$, $Y=X_0+\ldots+X_n$, $s=\sqrt{\Var(Y)}>0$, $S=\bover1sY$, maka untuk
setiap $\alpha\in[0,s/K]$

$$\Pr(|S|\ge\alpha)
\le 2\exp\bigl(\bover{-\alpha^2}{2(1+\bover{\alpha K}{2s})^2}\bigr)
\bumpeq 2e^{-\alpha^2/2}$$

\noindent jika $s\gg\alpha K$ ({\smc R\'enyi 70}, \S VII.4, Teorema 1).
Hasil yang kurang tajam dari jenis yang sama terdapat dalam 272Xl.

Sekarang saya mendaftarkan beberapa kasus baku ketika syarat Lindeberg
terpenuhi, sehingga kita dapat menerapkan teoremanya.
}%end of comment

\leader{274I}{Akibat} Misalkan $\sequencen{X_n}$ adalah barisan variabel
acak bernilai real yang saling bebas, semuanya dengan distribusi yang sama,
dan anggap bahwa ekspektasi bersamanya adalah $0$ serta varians bersamanya
berhingga dan tidak nol. Tuliskan $\sigma$ untuk nilai bersama
$\sqrt{\Var(X_n)}$, dan tetapkan

\Centerline{$S_n=\Bover1{\sigma\sqrt{n+1}}(X_0+\ldots+X_n)$}

\noindent untuk setiap $n\in\Bbb N$. Maka

\Centerline{$\lim_{n\to\infty}\Pr(S_n\le a)=\Phi(a)$}

\noindent secara seragam untuk $a\in\Bbb R$.

\proof{ Dalam bahasa 274G-274H, kita mempunyai $\sigma_n=\sigma$,
$s_n=\sigma\sqrt{n+1}$ dan $S_n=\Bover1{s_n}\sum_{i=0}^nX_i$.
Selain itu, jika $\nu$ adalah distribusi yang sama bagi semua $X_n$, maka

\Centerline{$\Expn(\psi_{\epsilon s_n}(X_n))
=\int_{\{x:|x|>\epsilon\sigma\sqrt{n+1}\}}x^2\nu(dx)\to 0$}

\noindent berdasarkan Teorema Konvergensi Terdominasi Lebesgue; sehingga

\Centerline{$\Bover1{s_n^2}\sum_{i=0}^n\Expn(\psi_{\epsilon s_n}(X_n))
\to 0$}

\noindent berdasarkan 273Ca. Jadi syarat Lindeberg terpenuhi dan 274G
memberikan hasilnya.
}%end of proof of 274I

\leader{274J}{Akibat} Misalkan $\sequencen{X_n}$ adalah barisan variabel
acak bernilai real yang saling bebas dengan ekspektasi nol, dan anggap bahwa
$\{X_n^2:n\in\Bbb N\}$ terintegralkan secara seragam dan bahwa

\Centerline{$\liminf_{n\to\infty}\Bover1{n+1}\sum_{i=0}^n\Var(X_i)>0$.}

\noindent Tetapkan

\Centerline{$s_n=\sqrt{\sumop_{i=0}^n\Var(X_i)}$,
\quad$S_n=\Bover1{s_n}(X_0+\ldots+X_n)$}
%\sumop needed because inside \sqrt{}

\noindent untuk $n\in\Bbb N$ yang besar. Maka

\Centerline{$\lim_{n\to\infty}\Pr(S_n\le a)=\Phi(a)$}

\noindent secara seragam untuk $a\in\Bbb R$.

\proof{ Syarat

\Centerline{$\liminf_{n\to\infty}\Bover1{n+1}\sum_{i=0}^n\Var(X_i)>0$}

\noindent berarti bahwa terdapat $c>0$, $n_0\in\Bbb N$ sedemikian sehingga
$s_n\ge c\sqrt{n+1}$ untuk setiap $n\ge n_0$. Misalkan ruang yang mendasarinya
adalah $(\Omega,\Sigma,\mu)$, dan ambil $\epsilon$, $\eta>0$. Dengan menulis
$\psi_{\delta}(x)=0$ untuk $|x|\le\delta$, dan $x^2$ untuk $|x|>\delta$,
seperti dalam 274F-274G, kita mempunyai

\Centerline{$\Expn(\psi_{\epsilon s_n}(X_i))
\le\Expn(\psi_{c\epsilon\sqrt{n+1}}(X_i))
=\int_{F(i,c\epsilon\sqrt{n+1})}X_i^2d\mu$}

\noindent untuk $n\ge n_0$, $i\le n$, dengan
$F(i,\gamma)=\{\omega:\omega\in\dom X_i,\,|X_i(\omega)|>\gamma\}$.
Karena $\{X_i^2:i\in\Bbb N\}$ terintegralkan secara seragam, terdapat
$\gamma\ge 0$ sedemikian sehingga $\int_{F(i,\gamma)}X_i^2d\mu\le\eta c^2$
untuk setiap $i\in\Bbb N$ (246I). Ambil $n_1\ge n_0$ sedemikian sehingga
$c\epsilon\sqrt{n_1+1}\ge\gamma$; maka untuk setiap $n\ge n_1$

\Centerline{$\Bover1{s_n^2}\sum_{i=0}^{n}\Expn(\psi_{\epsilon s_n}(X_i))
\le\Bover1{c^2(n+1)}\sum_{i=0}^n\eta c^2=\eta$.}

\noindent Karena $\epsilon$ dan $\eta$ sebarang, syarat-syarat 274G terpenuhi
dan hasilnya menyusul.
}%end of proof of 274J

\leader{274K}{Akibat} Misalkan $\sequencen{X_n}$ adalah barisan variabel
acak bernilai real yang saling bebas dengan ekspektasi nol, dan anggap bahwa

(i) terdapat suatu $\delta>0$ sedemikian sehingga
$\sup_{n\in\Bbb N}\Expn(|X_n|^{2+\delta})<\infty$,

(ii) $\liminf_{n\to\infty}\Bover1{n+1}\sum_{i=0}^n\Var(X_i)>0$.

\noindent Tetapkan $s_n=\sqrt{\sum_{i=0}^n\Var(X_i)}$ dan

\Centerline{$S_n=\Bover1{s_n}(X_0+\ldots+X_n)$}

\noindent untuk $n\in\Bbb N$ yang besar. Maka

\Centerline{$\lim_{n\to\infty}\Pr(S_n\le a)
=\Phi(a)$}

\noindent secara seragam untuk $a\in\Bbb R$.

\proof{ Intinya ialah bahwa $\{X_n^2:n\in\Bbb N\}$ terintegralkan secara
seragam. \Prf\ Tetapkan
$K=1+\sup_{n\in\Bbb N}\Expn(|X_n|^{2+\delta})$. Diberikan $\epsilon>0$,
tetapkan $M=(K/\epsilon)^{1/\delta}$. Maka
$(X_n^2-M)^+\le M^{-\delta}|X_n|^{2+\delta}$, sehingga

\Centerline{$\Expn(X_n^2-M)^+\le KM^{-\delta}=\epsilon$}

\noindent untuk setiap $n\in\Bbb N$. Karena $\epsilon$ sebarang,
$\{X_n^2:n\in\Bbb N\}$ terintegralkan secara seragam. \Qed

Dengan demikian syarat-syarat 274J terpenuhi dan kita memperoleh hasilnya.
}

\cmmnt{
\leader{274L}{Catatan (a)} Semua teorema dalam bagian ini ditujukan untuk
menemukan syarat-syarat ketika variabel acak $S$ `hampir' normal baku, dalam
arti $\Pr(S\le a)\bumpeq\Pr(Z\le a)$ secara seragam untuk $a\in\Bbb R$,
dengan $Z$ suatu (atau sebarang) variabel acak normal baku. Dalam semua kasus,
variabel acak $S$ dinormalisasi agar mempunyai ekspektasi $0$ dan varians $1$,
dan merupakan jumlah dari sejumlah besar variabel acak yang saling bebas.
(Dalam 274G dan 274I-274K dinyatakan secara eksplisit bahwa harus ada banyak
$X_i$, karena hasil-hasil tersebut merujuk pada limit ketika $n\to\infty$.
Hal ini tidak dinyatakan dengan kata-kata yang persis demikian dalam rumusan
teorema Lindeberg yang saya berikan, tetapi pembuktiannya memperjelas bahwa
$n(\delta+\delta^2)\ge 1$, sehingga tentu $n$ juga harus besar di sana.)

\header{274Lb}{\bf (b)} Saya tidak dapat meninggalkan bagian ini tanpa
mencatat bahwa bentuk definisi `hampir normal baku' dapat menyesatkan intuisi
Anda jika dicoba diterapkan pada distribusi lain. Jika kita mengambil $F$
sebagai fungsi distribusi $S$, sehingga $F(a)=\Pr(S\le a)$, saya mengatakan
bahwa $S$ `hampir normal baku' jika
$\sup_{a\in\Bbb R}|F(a)-\Phi(a)|$ kecil. Wajar untuk memandang ini sebagai
aproksimasi dalam suatu metrik, dengan menulis

\Centerline{$\tilde\rho(\nu,\nuprime)
=\sup_{a\in \Bbb R}|F_{\nu}(a)-F_{\nuprime}(a)|$}

\noindent untuk distribusi $\nu$, $\nuprime$ pada $\Bbb R$, dengan
$F_{\nu}(a)=\nu\ocint{-\infty,a}$. Dalam bentuk ini, teorema-teorema di atas
dapat dibaca sebagai pencarian syarat-syarat yang membuat
$\lim_{n\to\infty}\tilde\rho(\nu_{S_n},\mu_G)=0$. Namun intinya ialah bahwa
$\tilde\rho$ sebenarnya bukan metrik yang tepat digunakan. Metrik ini bekerja
di sini karena $\mu_G$ tak beratom. Tetapi misalkan, sebagai contoh, $\nu$
adalah ukuran Dirac pada $\Bbb R$ yang terpusat di $0$, dan $\nu_n$ adalah
distribusi variabel acak normal dengan ekspektasi $0$ dan varians
${1\over n}$, untuk setiap $n\ge 1$. Maka $F_{\nu}(0)=1$ dan
$F_{\nu_n}(0)={1\over 2}$, sehingga
$\tilde\rho(\nu_n,\nu)={1\over 2}$ untuk setiap $n\ge 1$. Namun, untuk
kebanyakan tujuan orang akan menganggap selisih antara $\nu_n$ dan $\nu$
kecil, dan tentunya $\nu$ adalah satu-satunya distribusi yang dapat disebut
secara wajar sebagai limit dari $\nu_n$.

\header{274Lc}{\bf (c)} Kesulitan-kesulitan di sini muncul dalam lebih dari
satu bentuk. Seorang statistikawan tidak akan menyukai gagasan bahwa $\nu_n$
dalam (b) di atas jauh dari $\nu$ (dan dari satu sama lain), dengan alasan
bahwa setiap pengukuran yang melibatkan variabel acak dengan distribusi ini
pasti mengalami galat, dan galat pengukuran yang kecil akan membuatnya tidak
dapat dibedakan. Seorang matematikawan murni, yang membayangkan kemungkinan
generalisasi hasil-hasil ini, tidak akan menyukai penekanan pada nilai
$\nu\ocint{-\infty,a}$, yang padanannya mungkin sulit ditemukan dalam ruang
yang lebih abstrak.

\header{274Ld}{\bf (d)} Pertimbangan-pertimbangan ini berpadu dan membawa
kita pada definisi yang agak berbeda untuk topologi pada ruang $P$ dari
distribusi probabilitas pada $\Bbb R$. Untuk setiap fungsi kontinu terbatas
$h:\Bbb R\to\Bbb R$ kita mempunyai pseudometrik
$\rho_h:P\times P\to\coint{0,\infty}$
yang didefinisikan dengan menulis

\Centerline{$\rho_h(\nu,\nuprime)=|\int h\,d\nu-\int h\,d\nuprime|$}

\noindent untuk semua $\nu$, $\nuprime\in P$. {\bf Topologi samar} pada $P$
adalah topologi yang dibangkitkan oleh pseudometrik $\rho_h$ (2A3F). Saya
tidak akan membahas sifat-sifatnya secara terperinci di sini (beberapa
diuraikan dalam 274Yc-274Yf %274Yc 274Yd 274Ye 274Yf
di bawah; lihat pula 285K-285L, 285S, dan 437J-437T
%437J 437K 437L 437M 437N (4{}37Q) 437P (4{}37T) 437R (4{}37O, 4{}37R)
%437S (4{}37P}) 437T (4{}37S)
dalam Jilid 4). Namun saya berpendapat bahwa cara yang tepat untuk memandang
hasil-hasil bab ini ialah dengan mengatakan bahwa (i) distribusi $\nu_S$
dekat dengan $\mu_G$ untuk topologi samar, (ii) himpunan
$\{\nu:\tilde\rho(\nu,\mu_G)<\epsilon\}$ terbuka untuk topologi tersebut,
dan itulah sebabnya $\tilde\rho(\nu_S,\mu_G)$ kecil.
}%end of comment

\leader{*274M}{}\cmmnt{ Saya menutup bagian ini dengan sepasang ketaksamaan
sederhana yang sering berguna ketika mempelajari variabel acak normal.

\medskip

\noindent}{\bf Lemma} (a)
$\int_x^{\infty}e^{-t^2/2}dt\le\Bover1xe^{-x^2/2}$ untuk setiap $x>0$.

(b) $\int_x^{\infty}e^{-t^2/2}dt\ge\Bover1{2x}e^{-x^2/2}$ untuk setiap
$x\ge 1$.

\proof{{\bf (a)}

$$\eqalign{\int_x^{\infty}e^{-t^2/2}dt
=\int_0^{\infty}e^{-(x+s)^2/2}ds
\le e^{-x^2/2}\int_0^{\infty}e^{-xs}ds
=\Bover1xe^{-x^2/2}.\cr}$$

\medskip

{\bf (b)} Tetapkan

\Centerline{$f(t)=e^{-t^2/2}-(1-x(t-x))e^{-x^2/2}$.}

\noindent Maka $f(x)=f'(x)=0$ dan
$f''(t)=(t^2-1)e^{-t^2/2}$ positif untuk $t\ge x$ (karena $x\ge 1$).
Dengan demikian $f(t)\ge 0$ untuk setiap $t\ge x$, dan
$\int_x^{x+1/x}f(t)dt\ge 0$. Namun ini tepat berarti bahwa

$$\eqalign{\int_x^{\infty}e^{-t^2/2}dt
\ge\int_x^{x+\bover1x}e^{-t^2/2}dt
\ge\int_x^{x+\bover1x}(1-x(t-x))e^{-x^2/2}dt
=\Bover1{2x}e^{-x^2/2},\cr}$$

\noindent sebagaimana diperlukan.
}%end of proof of 274M

\exercises{
\leader{274X}{Latihan dasar $\pmb{>}$(a)}
%\spheader 274Xa
Gunakan 272U untuk memberikan bukti alternatif bagi 274B.
%274B

\spheader 274Xb Misalkan $f:\Bbb R\to\Bbb R$ kontinu mutlak pada setiap
interval tertutup terbatas, dan bahwa
$\int_{-\infty}^{\infty}|f'(x)|e^{-ax^2}dx\penalty-100<\infty$ untuk setiap
$a>0$. Misalkan $X$ adalah variabel acak normal dengan ekspektasi nol.
Tunjukkan bahwa $\Expn(Xf(X))$ dan $\Expn(X^2)\Expn(f'(X))$ terdefinisi dan
sama.
%274B

\spheader 274Xc Buktikan 274D ketika $h''$ bersifat $M_3$-Lipschitz tetapi
tidak harus terdiferensialkan.
%274D

\spheader 274Xd Misalkan $\sequence{k}{m_k}$ adalah barisan yang meningkat
ketat dalam $\Bbb N$ sedemikian sehingga $m_0=0$ dan
$\lim_{k\to\infty}m_k/m_{k+1}=0$. Misalkan $\sequencen{X_n}$ adalah barisan
variabel acak yang saling bebas sedemikian sehingga
$\Pr(X_n=\sqrt{m_{k}})=\Pr(X_n=-\sqrt{m_{k}})=1/2m_{k}$,
$\Pr(X_n=0)=1-1/m_{k}$ kapan pun $m_{k-1}\le n<m_{k}$. Tunjukkan bahwa
Teorema Limit Pusat tidak berlaku untuk $\sequencen{X_n}$.
\Hint{dengan menetapkan $W_k=(X_0+\ldots+X_{m_k-1})/\sqrt{m_k}$, tunjukkan
bahwa $\Pr(W_k\in[\epsilon,1-\epsilon])\to 0$ untuk setiap $\epsilon>0$.}
%274G mt27bits

\spheader 274Xe Misalkan $\sequencen{X_n}$ adalah sebarang barisan variabel
acak yang saling bebas dan semuanya mempunyai distribusi yang sama; anggap
bahwa semuanya mempunyai varians berhingga $\sigma^2>0$, dan bahwa ekspektasi
bersamanya adalah $c$. Tetapkan
$S_n=c+\Bover1{\sqrt{n+1}}\sum_{i=0}^n(X_i-c)$ untuk setiap $n$, dan misalkan $Y$
adalah variabel acak normal dengan ekspektasi $c$ dan varians $\sigma^2$.
Tunjukkan bahwa $\lim_{n\to\infty}\Pr(S_n\le a)=\Pr(Y\le a)$
secara seragam untuk $a\in\Bbb R$.
%274I

%\sqheader 274Xf
\wheader{274Xf}{10}{4}{4}{48pt}$\pmb{>}${\bf (f)}
Tunjukkan bahwa untuk setiap $a\in\Bbb R$,

$$\lim_{n\to\infty}\bover{1}{2^n}
\sum_{r=0}^{\lfloor\bover{n}2+a\bover{\sqrt{n}}2\rfloor}
\bover{n!}{r!(n-r)!}
=\lim_{n\to\infty}\bover1{2^n}\#(\{I:I\subseteq
n,\,\#(I)\le\Bover{n}2+a\Bover{\sqrt{n}}2\})
=\Phi(a).$$
%274I

\spheader 274Xg Tunjukkan bahwa 274I adalah kasus khusus dari 274J.
%274J

\spheader 274Xh Misalkan $\sequencen{X_n}$ adalah barisan variabel acak
bernilai real yang saling bebas dengan ekspektasi nol. Tetapkan
$s_n=\sqrt{\sum_{i=0}^n\Var(X_i)}$ dan

\Centerline{$S_n=\Bover1{s_n}(X_0+\ldots+X_n)$}

\noindent untuk setiap $n\in\Bbb N$. Anggap bahwa terdapat suatu $\delta>0$
sedemikian sehingga

\Centerline{$\lim_{n\to\infty}\Bover1{s_n^{2+\delta}}
\sum_{i=0}^n\Expn(|X_i|^{2+\delta})=0$.}

\noindent Tunjukkan bahwa $\lim_{n\to\infty}\Pr(S_n\le a)=\Phi(a)$ secara
seragam untuk $a\in\Bbb R$. (Ini adalah suatu bentuk {\bf teorema limit pusat
Liapounoff}; lihat {\smc Liapounoff 1901}.)
%274K

\spheader 274Xi Misalkan $P$ adalah himpunan ukuran probabilitas Radon pada
$\Bbb R$. Misalkan $\nu_0\in P$, $a\in\Bbb R$. Tunjukkan bahwa pemetaan
$\nu\mapsto\nu\ocint{-\infty,a}:P\to[0,1]$ kontinu di $\nu_0$ untuk
topologi samar pada $P$ jika dan hanya jika $\nu_0\{a\}=0$.
%274L

\spheader 274Xj Misalkan $\sequencen{X_n}$ adalah barisan variabel acak
yang saling bebas dan berdistribusi identik dengan varians berhingga tak nol.
Misalkan $\sequencen{t_n}$ adalah barisan dalam $\Bbb R$ sedemikian sehingga
$\sum_{n=0}^{\infty}t_n^2=\infty$. Tunjukkan bahwa
$\sum_{n=0}^{\infty}t_nX_n$ tidak terdefinisi atau tak berhingga hampir di
mana-mana. \Hint{Tangani dahulu kasus ketika $\sequencen{t_n}$ tidak
konvergen ke $0$. Jika tidak demikian, gunakan 274G untuk menunjukkan bahwa,
untuk setiap $n\in\Bbb N$,
$\lim_{m\to\infty}\Pr(|\sum_{i=n}^mt_iX_i|\ge 1)\ge\bover12$. Lihat pula
276Xd.}
%274G
%query:  do we really need `finite variance'?

\spheader 274Xk Misalkan $\sequencen{X_n}$ adalah barisan variabel acak
bernilai real yang saling bebas dengan ekspektasi nol. Anggap bahwa $M\ge 0$
sedemikian sehingga $|X_n|\le M$ hampir di mana-mana untuk setiap $n$, dan
bahwa $\sum_{n=0}^{\infty}\Var(X_n)=\infty$. Tetapkan
$s_n=\sqrt{\sum_{i=0}^n\Var(X_i)}$ untuk setiap $n$, dan
$S_n=\Bover1{s_n}\sum_{i=0}^nX_i$ ketika $s_n>0$. Tunjukkan bahwa
$\lim_{n\to\infty}\Pr(S_n\le a)=\Phi(a)$ untuk setiap $a\in\Bbb R$.
%274G
%out of order query

\leader{274Y}{Latihan lanjutan (a)}
%\spheader 274Ya
({\smc Steele 86})
Misalkan $X_0,\ldots,X_n,Y_0,\ldots,Y_n$ adalah variabel acak yang saling
bebas sedemikian sehingga, untuk setiap $i\le n$,
$X_i$ dan $Y_i$ mempunyai distribusi yang sama.
Misalkan $h:\BbbR^{n+1}\to\Bbb R$ adalah fungsi terukur Borel, dan tetapkan
$Z=h(X_0,\ldots,X_n)$,
$Z_i=h(X_0,\ldots,X_{i-1},Y_i,X_{i+1},\ldots,X_n)$ untuk setiap $i$
(dengan $Z_0=h(Y_0,X_1,\ldots,X_n)$ dan
$Z_n=h(X_0,\ldots,X_{n-1},Y_n)$, tentu saja). Anggap bahwa $Z$ mempunyai
ekspektasi berhingga. Tunjukkan bahwa
$\Var(Z)\le\bover12\sum_{i=0}^n\Expn(Z_i-Z)^2$.
%274C mt27bits

\spheader 274Yb\dvAnew{2010} Tunjukkan bahwa untuk setiap $\epsilon>0$
terdapat fungsi mulus $h:\Bbb R\to[0,1]$ sedemikian sehingga
$\chi\ocint{-\infty,-\epsilon}\le h\le\chi\ocint{-\infty,\epsilon}$.
%274E

\spheader 274Yc\dvAformerly{2{}74Ya}
Tuliskan $P$ untuk himpunan ukuran probabilitas Radon pada $\Bbb R$. Untuk
$\nu$, $\nuprime\in P$ tetapkan

$$\eqalign{\rho(\nu,\nuprime)
=\inf\{\epsilon:\epsilon\ge 0,\,\nu\ocint{-\infty,a-\epsilon}-\epsilon
\le\nuprime\ocint{-\infty,a}
&\le\nu\ocint{-\infty,a+\epsilon}+\epsilon\cr
&\quad\quad\text{ untuk setiap }a\in\Bbb R\}.\cr}$$

\noindent Tunjukkan bahwa $\rho$ adalah metrik pada $P$ dan mendefinisikan
topologi samar pada $P$. ($\rho$ disebut {\bf metrik L\'evy}.)
%274L

\spheader 274Yd Tuliskan $P$ untuk himpunan ukuran probabilitas Radon pada
$\Bbb R$, dan misalkan $\tilde\rho$ adalah metrik pada $P$ yang didefinisikan
dalam 274Lb. Tunjukkan bahwa jika $\nu\in P$ tak beratom dan $\epsilon>0$,
maka $\{\nuprime:\nuprime\in P,\,\tilde\rho(\nuprime,\nu)<\epsilon\}$ terbuka
untuk topologi samar pada $P$.
%274L

\spheader 274Ye Misalkan $\sequencen{S_n}$ adalah barisan variabel acak
bernilai real, dan $Z$ variabel acak normal baku. Tunjukkan bahwa pernyataan
berikut ekuivalen:

\quad(i) $\mu_G=\lim_{n\to\infty}\nu_{S_n}$ untuk topologi samar, dengan
$\nu_{S_n}$ menyatakan distribusi $S_n$;

\quad(ii) $\Expn(h(Z))=\lim_{n\to\infty}\Expn(h(S_n))$ untuk setiap fungsi
kontinu terbatas $h:\Bbb R\to\Bbb R$;

\quad(iii) $\Expn(h(Z))=\lim_{n\to\infty}\Expn(h(S_n))$ untuk setiap fungsi
terbatas $h:\Bbb R\to\Bbb R$ sedemikian sehingga
($\alpha$) $h$ mempunyai turunan kontinu dari semua orde
($\beta$) $\{x:h(x)\ne 0\}$ terbatas;

\quad(iv) $\lim_{n\to\infty}\Pr(S_n\le a)=\Phi(a)$ untuk setiap
$a\in\Bbb R$;

\quad(v) $\lim_{n\to\infty}\Pr(S_n\le a)=\Phi(a)$ secara seragam untuk
$a\in\Bbb R$;

\quad(vi) $\{a:\lim_{n\to\infty}\Pr(S_n\le a)=\Phi(a)\}$ rapat dalam
$\Bbb R$.

\noindent (Lihat pula 285L.)
%274L

\spheader 274Yf Misalkan $(\Omega,\Sigma,\mu)$ adalah ruang probabilitas,
dan $P$ himpunan ukuran probabilitas Radon pada $\Bbb R$. Tunjukkan bahwa
$X\mapsto\nu_X:\eusm L^0(\mu)\to P$ kontinu untuk topologi konvergensi dalam
ukuran pada $\eusm L^0(\mu)$ dan topologi samar pada $P$.
%274L

\spheader 274Yg Misalkan $\sequencen{X_n}$ adalah barisan variabel acak
bernilai real yang saling bebas. Anggap bahwa terdapat $M\ge 0$ sedemikian
sehingga $|X_n|\le M$ hampir di mana-mana untuk setiap $n\in\Bbb N$, dan
bahwa $\sum_{n=0}^{\infty}X_n$ terdefinisi, sebagai bilangan real, hampir di
mana-mana. Tunjukkan bahwa $\sum_{n=0}^{\infty}\Var(X_n)<\infty$.
%274Xk 274G 274Yf

}%end of exercises

\endnotes{
\Notesheader{274} Selama lebih dari dua ratus tahun Teorema Limit Pusat telah
menjadi salah satu kejayaan matematika, dan tidak ada cabang matematika atau
sains yang akan tetap sama tanpanya. Saya kira teorema ini adalah satu teorema
terpenting dalam teori probabilitas; dan saya mengamati bahwa pembuktiannya
hampir tidak menggunakan teori ukuran. Memang, saya telah membungkus argumen
di atas dalam bahasa ukuran dan integrasi. Namun jika kita melihat hakikatnya,
unsur-unsur penting pembuktian tersebut adalah

\inset{(i) kombinasi linear variabel acak normal yang saling bebas adalah
normal (274Ae, 274B);}

\inset{(ii) jika $U$, $V$, $W$ adalah variabel acak yang saling bebas, dan
$h$ fungsi kontinu terbatas, maka
$|\Expn(h(U,V,W))|\le\sup_{t\in\Bbb R}|\Expn(h(U,V,t))|$ (274C);}

\inset{(iii) jika $(X_0,\ldots,X_n)$ adalah variabel acak yang saling bebas,
maka kita dapat menemukan variabel acak yang saling bebas
$(X_0',\ldots,X_n',Z_0,\ldots,Z_n)$ sedemikian sehingga $Z_j$ normal baku dan
$X_j'$ mempunyai distribusi yang sama dengan $X_j$, untuk setiap $j$ (274F).}

\noindent Bagian selebihnya dari argumen terdiri atas kalkulus elementer,
taksiran cermat, dan beberapa sifat paling mendasar dari ekspektasi dan
kebebasan. Sekarang (ii) dan (iii) dibenarkan di atas dengan menggunakan
teorema Fubini, tetapi tentu saja keduanya termasuk dalam daftar intuisi
probabilistik yang mendahului pengidentifikasian probabilitas dengan fungsional
aditif terhitung. Seandainya keduanya menimbulkan kesulitan yang tidak dapat
diatasi, hal itu akan menjadi argumen kuat yang menentang model probabilitas
yang kita gunakan, tetapi tidak akan memengaruhi Teorema Limit Pusat.
Sesungguhnya (i) tampaknya merupakan tempat kita benar-benar memerlukan model
matematis bagi konsep `distribusi', dan semua perhitungan yang relevan dapat
dilakukan dalam pengertian integral Riemann pada bidang, tanpa menyebut
aditivitas terhitung. Jadi, meskipun saya senang dan bangga telah menuliskan
suatu versi dari gagasan-gagasan indah ini, saya harus mengakui bahwa semuanya
tidak bergantung secara esensial pada bagian lain risalah ini.

Dalam \S285 saya akan mendeskripsikan pendekatan yang sama sekali berbeda
terhadap teorema ini, dengan menggunakan perangkat yang jauh lebih canggih;
tetapi sekali lagi, barangkali dengan lebih tersembunyi, relevansi teori ukuran
bukan terletak pada teorema itu sendiri, melainkan pada bayangan kita tentang
apa yang dimaksud dengan distribusi sebarang. Sebab di sinilah saya memang
mengajukan klaim bagi bidang saya. Karakterisasi fungsi distribusi sebagai
fungsi monoton sebarang, kontinu dari kanan, dan dengan limit yang tepat pada
$\pm\infty$ (271Xb), bersama dengan analisis fungsi monoton dalam \S226,
memberi kita kesempatan untuk membentuk gambaran mental tentang kelas objek
yang tepat bagi penerapan hasil-hasil seperti Teorema Limit Pusat.

Teorema 274F adalah modifikasi kecil dari Teorema 3 dalam
{\smc Lindeberg 1922}. Seperti aslinya, teorema ini menekankan apa yang saya
yakini penting bagi semua teorema limit dalam bab ini: teorema-teorema tersebut
sebaiknya didasarkan pada pemahaman yang tepat mengenai barisan berhingga
variabel acak. Syarat Lindeberg merupakan puncak dari pencarian panjang akan
syarat paling umum yang membuat Teorema Limit Pusat berlaku. Saya menawarkan
suatu versi teorema Laplace (274Xf) sebagai titik awal, dan syarat Liapounoff
(274Xh) sebagai contoh salah satu tahap perantara. Tentu saja akibat 274I,
274J, 274K, dan 274Xe adalah bentuk yang hendak kita terapkan jika dapat
memilih. Terdapat kesejajaran yang menarik, tetapi sejauh yang saya ketahui
murni kebetulan, antara 273H/274K dan 273I/274Xe. Sebagai contoh barisan
variabel acak $\sequencen{X_n}$ yang saling bebas, semuanya dengan ekspektasi
nol dan varians $1$, tetapi yang {\it tidak} memenuhi Teorema Limit Pusat, saya
menawarkan 274Xd.
}%end of notes

\discrpage
